%% ============================================================
%%  Chapter 4: Experimental Evaluation
%% ============================================================
\chapter{Experimental Evaluation}
\label{chap:experiments}

% GUIDANCE: 15--20 pages. This is your results chapter.
% Write the setup sections NOW (they don't need results).
% Leave \todo placeholders for figures and tables.
% Fill in numbers once experiments finish.


% ── 4.1 ─────────────────────────────────────────────────────────
\section{Experimental Setup}
\label{sec:exp_setup}


\subsection{Environment Configuration}
\label{subsec:exp_environment}

\todo{Describe the simulation environment used for all experiments:
  30$\times$30 grid, 6 agents, 8 induct stations, 38 eject stations.
  YAML config file parameters. Timestep definition.}


\subsection{Parameters Swept}
\label{subsec:exp_parameters}

All experiments sweep the following parameters:

\begin{itemize}
  \item \textbf{Communication range} $\comm \in \{3, 5, 8, 13, 20, 45\}$
        grid units, ranging from fully isolated to fully connected.
  \item \textbf{Tasks per induct station} $\in \{2, 5, 10\}$,
        yielding 16, 40, and 80 total tasks respectively.
  \item \textbf{Rerun interval} (Dynamic GCBBA only)
        $\in \{10, 25, 50, 100, 200\}$ timesteps, with $ri = 50$
        as the canonical configuration.
\end{itemize}

Each configuration is repeated across 5 random seeds.
The maximum timestep budget is 800 per run.

\todo{Add table summarising all parameter combinations and
total run counts per mode (quick/medium/full).}


\subsection{Metrics}
\label{subsec:exp_metrics}

\todo{Define each metric precisely:
  \textbf{Makespan}: timesteps to complete all tasks ($-1$ if DNF);
  \textbf{Completion rate}: fraction of seeds where all tasks complete;
  \textbf{Sub-optimality ratio}: method makespan / SGA makespan
    (only computed where both methods complete);
  \textbf{Distinct stuck events}: number of agent transitions
    into the stuck state (NOT cumulative stuck timesteps);
  \textbf{Vertex collisions}: post-hoc count of two agents at
    the same cell at the same timestep;
  \textbf{Task balance std}: std dev of tasks completed per agent;
  \textbf{Average idle ratio}: fraction of timesteps each agent
    spends not moving.}


\subsection{Baselines}
\label{subsec:exp_baselines}

\todo{One paragraph each on SGA and CBBA as implemented.
  Reiterate that all methods share the same collision avoidance
  pipeline -- this is the key design decision for fair comparison.
  Note the SGA caveat: on disconnected graphs, SGA runs per
  connected component independently (generous upper bound).}


% ── 4.2 ─────────────────────────────────────────────────────────
\section{Communication Robustness}
\label{sec:exp_comm_robustness}

% THE KEY RESULT. Dynamic GCBBA completes tasks where static fails.

\subsection{Task Completion Rate}
\label{subsec:exp_completion_rate}

\todo{Write analysis of completion rate results.
  Key finding: Dynamic GCBBA achieves 100\% completion at $\comm = 13$
  (barely connected) and disconnected ranges where Static GCBBA and
  CBBA achieve 0\%. Reference Figure~\ref{fig:completion_rate}.}

\begin{figure}[htbp]
  \centering
  % \includegraphics[width=\textwidth]{completion_rate}
  \todo{Insert completion\_rate.png}
  \caption{Task completion rate by communication range, aggregated
           across all task loads and seeds. The shaded region
           indicates disconnected graph configurations
           ($\comm < 13$).}
  \label{fig:completion_rate}
\end{figure}

\begin{figure}[htbp]
  \centering
  % \includegraphics[width=\textwidth]{completion_rate_by_tpi}
  \todo{Insert completion\_rate\_by\_tpi.png}
  \caption{Task completion rate broken down by task load.}
  \label{fig:completion_rate_by_tpi}
\end{figure}


\subsection{Makespan Under Varying Connectivity}
\label{subsec:exp_makespan_comm}

\todo{Analyse makespan vs communication range.
  For connected graphs ($\comm \geq 13$), all methods complete
  tasks -- focus on efficiency differences.
  For disconnected graphs, only Dynamic GCBBA completes --
  discuss why (event-driven replanning allows agents to
  reassign when connectivity changes as robots move).
  Reference Figure~\ref{fig:makespan_vs_comm}.}

\begin{figure}[htbp]
  \centering
  % \includegraphics[width=\textwidth]{makespan_vs_comm_range}
  \todo{Insert makespan\_vs\_comm\_range.png}
  \caption{Makespan vs.\ communication range for all methods
           across three task loads.}
  \label{fig:makespan_vs_comm}
\end{figure}


% ── 4.3 ─────────────────────────────────────────────────────────
\section{Optimality Analysis}
\label{sec:exp_optimality}

% Second key result: GCBBA stays near SGA on connected graphs.

\subsection{Sub-Optimality Ratio vs.\ SGA}
\label{subsec:exp_suboptimality}

\todo{Present sub-optimality ratio results.
  Key message: on connected graphs, Dynamic GCBBA achieves
  ratio $\approx 1.0$--$1.X$, well within the theoretical
  bound of 2.0. Static GCBBA shows higher ratio due to
  inability to reallocate as tasks complete.
  Reference Figure~\ref{fig:optimality_ratio}.}

\begin{figure}[htbp]
  \centering
  % \includegraphics[width=\textwidth]{optimality_ratio}
  \todo{Insert optimality\_ratio.png. Include the theoretical
  bound line at 2.0 and the SGA reference line at 1.0.}
  \caption{Sub-optimality ratio (method makespan / SGA makespan)
           vs.\ communication range. Only runs where both the
           method and SGA completed are included.
           The dashed line at 2.0 indicates the theoretical
           worst-case guarantee for GCBBA on connected graphs.}
  \label{fig:optimality_ratio}
\end{figure}


\subsection{Makespan Gap Relative to SGA}
\label{subsec:exp_makespan_gap}

\todo{Analyse improvement\_ratio.png.
  Show that on fully connected graphs, Dynamic GCBBA is within
  X\% of SGA. Note that on disconnected graphs, SGA is not a
  fair reference (it runs per component).}

\begin{figure}[htbp]
  \centering
  % \includegraphics[width=\textwidth]{improvement_ratio}
  \todo{Insert improvement\_ratio.png}
  \caption{Makespan gap relative to SGA, expressed as a
           percentage. Positive values indicate slower
           performance than SGA. Only runs where both the
           method and SGA completed are included.}
  \label{fig:improvement_ratio}
\end{figure}


% ── 4.4 ─────────────────────────────────────────────────────────
\section{Rerun Interval Sensitivity}
\label{sec:exp_sensitivity}

\todo{Present the rerun interval sweep results.
  Key message: performance is not critically sensitive to the
  chosen value of $ri$, validating the $ri = 50$ design choice.
  Very low $ri$ (10) may degrade due to computational overhead
  of frequent replanning. Very high $ri$ (200) degrades
  completion rate at lower comm ranges.
  Reference Figure~\ref{fig:ri_sweep}.}

\begin{figure}[htbp]
  \centering
  % \includegraphics[width=\textwidth]{rerun_interval_sweep}
  \todo{Insert rerun\_interval\_sweep.png}
  \caption{Makespan (top) and completion rate (bottom) for
           Dynamic GCBBA across rerun interval values.
           The canonical $ri = 50$ is shown with a solid line.}
  \label{fig:ri_sweep}
\end{figure}


% ── 4.5 ─────────────────────────────────────────────────────────
\section{Safety: Collisions and Deadlocks}
\label{sec:exp_safety}

\todo{Present collision and deadlock results.
  Expected finding: vertex collisions should be zero or near-zero
  (the collision avoidance pipeline works). If non-zero, explain
  under what conditions they occur.
  Stuck events should decrease with higher comm range / lower
  task load. Reference Figure~\ref{fig:collision_deadlock}.
  Clearly state that stuck events are transition-based
  (one event per agent entry into stuck state, not per timestep).}

\begin{figure}[htbp]
  \centering
  % \includegraphics[width=\textwidth]{collision_deadlock}
  \todo{Insert collision\_deadlock.png}
  \caption{Post-hoc vertex collision count (left) and distinct
           agent stuck events (right) by communication range.}
  \label{fig:collision_deadlock}
\end{figure}


% ── 4.6 ─────────────────────────────────────────────────────────
\section{Computational Overhead}
\label{sec:exp_timing}

\todo{Present GCBBA timing results.
  Key question: does the extra overhead of dynamic replanning
  justify the completion rate gains?
  Show avg allocation time per call vs task count.
  Show number of allocation reruns vs comm range.
  Dynamic GCBBA does more reallocation calls -- is the wall-clock
  cost acceptable for real-time use?
  Reference Figure~\ref{fig:gcbba_timing} and
  Figure~\ref{fig:trigger_breakdown}.}

\begin{figure}[htbp]
  \centering
  % \includegraphics[width=\textwidth]{gcbba_timing}
  \todo{Insert gcbba\_timing.png}
  \caption{Average allocation computation time per call (left)
           and number of allocation reruns per run (right),
           grouped by task load.}
  \label{fig:gcbba_timing}
\end{figure}

\begin{figure}[htbp]
  \centering
  % \includegraphics[width=\textwidth]{trigger_breakdown}
  \todo{Insert trigger\_breakdown.png}
  \caption{GCBBA rerun trigger breakdown for Dynamic GCBBA
           ($ri = 50$): batch-triggered (task completions)
           vs.\ interval-triggered reruns.}
  \label{fig:trigger_breakdown}
\end{figure}


% ── 4.7 ─────────────────────────────────────────────────────────
\section{Agent Utilisation and Load Balance}
\label{sec:exp_utilisation}

\todo{Present agent utilisation results.
  Key question: does GCBBA distribute tasks fairly across agents?
  Compare task\_balance\_std across methods.
  Dynamic GCBBA should show better balance (replanning corrects
  early imbalances as tasks complete).
  Reference Figure~\ref{fig:utilisation}.}

\begin{figure}[htbp]
  \centering
  % \includegraphics[width=\textwidth]{agent_utilization}
  \todo{Insert agent\_utilization.png}
  \caption{Agent idle ratio (left) and task balance standard
           deviation (right) by task load. Lower balance std
           indicates more even task distribution.}
  \label{fig:utilisation}
\end{figure}


% ── 4.8 ─────────────────────────────────────────────────────────
\section{Throughput Analysis}
\label{sec:exp_throughput}

\todo{Present throughput curves (cumulative tasks completed
  over time) for the most informative configurations:
  lowest comm range (most stress) and highest (baseline).
  Show the difference in curve shape between Static and Dynamic
  GCBBA at low comm ranges -- Static plateaus early, Dynamic
  continues completing tasks.
  Reference Figure~\ref{fig:throughput}.}

\begin{figure}[htbp]
  \centering
  % \includegraphics[width=\textwidth]{throughput_curves}
  \todo{Insert throughput\_curves.png}
  \caption{Cumulative task completion over time at the lowest
           (left) and highest (right) communication ranges,
           for maximum task load.}
  \label{fig:throughput}
\end{figure}


% ── 4.9 ─────────────────────────────────────────────────────────
\section{Summary of Results}
\label{sec:exp_summary}

\todo{Write a 1-page synthesis. Use the results table
  (Table~\ref{tab:results}) as the anchor. Structured around
  the three research questions from Chapter 1:
  RQ1 (completion under disconnected graphs): answered by
    Section~\ref{sec:exp_comm_robustness};
  RQ2 (near-optimality on connected graphs): answered by
    Section~\ref{sec:exp_optimality};
  RQ3 (replanning frequency): answered by
    Section~\ref{sec:exp_sensitivity}.}

\begin{table}[htbp]
  \centering
  \todo{Insert results\_table.tex generated by plot\_results.py.
  Use \textbackslash input\{results\_table\} once the file is
  generated.}
  \caption{Summary of experimental results across all methods,
           task loads, and communication ranges.}
  \label{tab:results}
\end{table}
